\begin{proposition}[An independent groupoid-suspension replacement for C.15]
\label{F:C15-groupoid-suspension}
Let $G$ be a nontrivial group.  There is a marked ordinary category
with a globally invariant diagram-localization presheaf on all small
tests whose infinity localization has a mapping space $BG$.
\end{proposition}
\begin{proof}
Let $\mathcal M$ have objects $a,x,b$, with $\operatorname{End}(x)=G$,
identity endomorphisms at $a,b$, and unique further arrows
$u:x\to a$ and $v:x\to b$.  Thus $ug=u$ and $vg=v$.
Mark $u$, all elements of $G$, and the identities, but not $v$.

Fix a path type $\tau$ with vertices $0,\ldots,L$.  Its negative
horizontal arrows must be marked; its interior vertical arrows are
arbitrary, and the endpoint vertical arrows are identities.
Write $K_\tau(a,b)$ for this path category.  In any such path an
entry $a\to x$ uses a negative $u$, an exit $x\to a$ uses a
positive $u$, and an exit $x\to b$ uses a positive $v$.
Once $b$ occurs it persists, since a negative $v$ is forbidden.

First normalize the horizontal group labels.  On each maximal
consecutive $x$ interval, set the gauge at its first vertex to $1$
and transport along successive horizontal arrows, using inverses
where their orientations require them.  Conjugating vertical
components by these gauges defines a normalization functor and a
natural isomorphism from the identity to it.  Its image is the full
category of paths whose horizontal group labels are all $1$.

For a normalized path $p$, let $I(p)$ be its last nonempty interval
of $x$ vertices.  Replace every earlier $x$ interval by $a$ and
denote the resulting path by $T(p)$.  For any ladder $p\to q$ one has
\[
 I(q)\subseteq I(p).
\]
Indeed source $a$ vertices must remain $a$, and source $b$ vertices
must remain $b$.  The last target $x$ interval lies after the last
source $a$ vertex: otherwise the target would already have become
$b$ there.  It also lies before the first source $b$ vertex.  The
source vertices strictly between these two vertices are exactly
$I(p)$.  On the retained target interval the normalized vertical
$G$ components are equal.  Use that component for $T(p)\to T(q)$;
all other components are forced identities, $u$'s or $v$'s.
Containment proves compatibility with composition at a common
vertex.  The other naturality equations are $ug=u$ and $vg=v$.
Thus $T$ is a functor, and the components $u$ on the discarded
intervals give a natural transformation $\mathrm{id}\Rightarrow T$.

The image of $T$ is the category $\mathbf BG\times P_\tau$, where
$P_\tau$ consists of admissible nonempty $x$ intervals, ordered by
reverse inclusion.  Admissibility means that the entry edge is
negative and the exit edge is positive.  A ladder between two such
paths is exactly one freely chosen element of $G$ on the smaller
interval, which proves the product description including composition.
When the stage is nonempty, $P_\tau$ has initial object
\[
 I_{\max}=[\text{first negative edge},
                 \text{last positive edge}-1].
\]
It contains every admissible interval.  Projection to $\mathbf BG$
and inclusion at this initial object therefore give a specified
categorical homotopy equivalence.  Combined with the two preceding
normalizations, this proves
\[
 K_\tau(a,b)\simeq\mathbf BG\quad\text{in }\mathrm{Cat}_{\mathrm{hot}}.
\]
These finite zigzags of natural transformations remain such after
$\operatorname{Fun}(Y,-)$ for every small $Y$.

The specified comparisons also control transitions.  On the canonical
$\mathbf BG$ inclusion all horizontal group labels are identities.
Path pullbacks, subdivisions, identity insertions and permitted
compositions retain those identities and retain a vertical element
$g$ at every surviving $x$ vertex.  Subsequent normalization and
last-interval contraction therefore return that same $g$.  All
actual transitions induce the identity of $\mathbf BG$ in
$\mathrm{Cat}_{\mathrm{hot}}$.  The nonempty stages are cofinal.
Using the established ordinary diagram-localization path formula,
\[
 H_{a,b}(Y)=\colim_\tau\pi_0\operatorname{Fun}(Y,K_\tau(a,b))
          =\pi_0\operatorname{Fun}(Y,\mathbf BG)
\]
for every small $Y$.  Functors to a groupoid factor through
$\Pi_1Y$, so this functor inverts every nerve weak equivalence,
without restricting tests to finite categories.

This verifies global invariance for all endpoint pairs as well.
Every $a$-to-$a$ stage has the constant $a$ path as terminal object;
$b$ cannot occur and the needed components are the unique $u$'s.
Every $b$-to-$b$ stage is a point.  There are no $b$-to-$a$ paths,
giving the empty value for nonempty $Y$ and a singleton for empty
$Y$.  Finally the marked natural transformation $u_Y:x_Y\to a_Y$
is inverted in each diagram localization, identifying every
remaining endpoint pair naturally with one of these pairs.

For the infinity-localization computation no general theorem about
arbitrary vertical ladders is needed.  The nerve $N\mathcal M$ is
the union of two right cones on $N\mathbf BG$, with cone vertices
$a,b$ and intersection $N\mathbf BG$.  Both inclusions are
monomorphisms, so this cofibrant span is a homotopy pushout in the
Joyal model.  A functor to an infinity category $\mathcal D$ is
therefore a functor $F:BG\to\mathcal D$, objects $A,B$, and cones
$\alpha:F\to\operatorname{const}A$ and
$\beta:F\to\operatorname{const}B$.  Inverting $u$ means that
$\alpha$ is a pointwise equivalence.  For fixed $A$ the choices
of $F$ with this specified equivalence are contractible.  The
remaining datum has space
\[
 \operatorname{Map}\bigl(BG,\operatorname{Map}_{\mathcal D}(A,B)\bigr).
\]
This is precisely the universal property of the two-object infinity
category with identity endomorphisms, no maps $b\to a$, and
$\operatorname{Map}(a,b)=BG$.  Its mapping component has fundamental
group $G\ne1$ and hence is not discrete.  This proves the desired
example independently of the unavailable original C.15 wording.
For the finite instance $G=C_2$, the program
\path{code/f_groupoid_suspension_table.py} supplies all six arrows,
all defined products, every identity and associativity check, and
closure of the five-arrow marking. The general stage argument above
is independent of any finite bound on the test category.
An independent stronger certificate,
\path{code/b_F_C15_groupoid_suspension.py}, constructs all sixteen
nonempty $a$-to-$b$ stages of length at most four, their actual
normalization and last-interval functors, and the specified natural
transformations. It checks all 25\,408 typed associativity triples,
746 natural-transformation equations, and 52 identity-insertion or
same-sign-composition transition functors. Every induced map on
$C_2$ is the identity. Complete tables and comparisons are stored in
\path{code/outputs/b_F_C15_certificate.json.gz}; the general proof,
rather than this finite bound, establishes the all-small-test claim.
\end{proof}
